% Five verified computational primitives — pre-submission LaTeX source
% Compile: pdflatex main && bibtex main && pdflatex main && pdflatex main
% Or: latexmk -pdf main.tex

\documentclass[11pt,reqno]{amsart}

\usepackage[utf8]{inputenc}
\usepackage[T1]{fontenc}
\usepackage{amsmath,amsthm,amsfonts,amssymb}
\usepackage{mathrsfs}
\usepackage{microtype}
\usepackage[margin=1in]{geometry}
\usepackage{enumitem}
\usepackage{booktabs}
\usepackage{array}
\usepackage{verbatim}
\usepackage{xcolor}
\usepackage{hyperref}
\hypersetup{
    colorlinks=true,
    linkcolor=blue!50!black,
    citecolor=blue!50!black,
    urlcolor=blue!50!black
}
\usepackage{cleveref}

% Theorem-like environments
\newtheorem{theorem}{Theorem}[section]
\newtheorem{proposition}[theorem]{Proposition}
\newtheorem{lemma}[theorem]{Lemma}
\newtheorem{corollary}[theorem]{Corollary}

\theoremstyle{definition}
\newtheorem{remark}[theorem]{Remark}

% Operators
\DeclareMathOperator{\erf}{erf}
\DeclareMathOperator{\erfc}{erfc}
\DeclareMathOperator{\sech}{sech}
\DeclareMathOperator{\csch}{csch}
\DeclareMathOperator{\Ai}{Ai}

% Custom
\newcommand{\dps}{\mathrm{dps}}
\newcommand{\rkc}{\texttt{rkc}}

\title{Five verified computational primitives via cross-domain transfer
of a Ramanujan-style operator basis: implementation, benchmarks, prior art}

\author{Jouni Helminen}
\address{Independent researcher}
\email{jouni@helminen.co}

\date{May 25, 2026}

\subjclass[2020]{Primary 65D20, 65Y20; Secondary 33C10, 33E12, 68T07, 82B20}

\keywords{Computational primitives; state-space attention; annular Green
function; Schottky--Klein prime function; harmonic-chain propagator;
Airy function; Tricomi confluent hypergeometric function; Madelung
constant; Ewald summation; operator basis; cross-domain transfer}

\begin{document}

\begin{abstract}
We present five computational primitives obtained by applying a small
operator basis (developed in a separate Ramanujan-mining project) to
problem classes outside its native number-theoretic domain. Each
primitive replaces a slower iterative or large-$N$ evaluator with a
closed-form or sparse-sum expression: (a)~finite-pole semiseparable
attention with $O(N \cdot R \cdot d_v)$ recurrence in place of $O(N^2)$
softmax direct form; (b)~annular Poisson Green function via one-image
$\theta_1$-quotient in place of $M$-mode Fourier series at thin-annulus
thickness $L \ll 1$; (c)~periodic 1D harmonic-chain port-Hamiltonian
propagator as a single Bessel-image sum in place of explicit
time-stepping; (d)~Airy/Bessel asymptotic prefactor via two-sector
Tricomi-$U$ closed form in place of Borel--Pad\'e quadrature;
(e)~CsCl-lattice signed Madelung constant via $42$ erfc/exp shells in
place of a $1{,}201$-term 2D csch/sech sum. All five are implemented in
the \rkc\ package and an accompanying analysis script, verified against
converged baselines at $30+$ decimal digits, benchmarked, and
accompanied by explicit prior-art attribution. We frame this as a
\emph{verified-implementation tool paper}, not a novelty claim:
(a)~is the multi-pole extension of RetNet/Mamba-SSD; (b)~is the
thin-annulus packaging of Crowdy--Marshall $\theta$-quotient Green
functions; (c)~is Schr\"odinger's 1914 chain extended via standard
image-method periodicity per Trovato \emph{et al.}~(2016); (d)~re-derives
a DLMF-adjacent representation with specific $(1/6, 1/3)$ parameters
less documented in the canonical references; (e)~re-packages classical
Ewald/Poisson machinery (Borwein--Borwein--Taylor; Crandall--Buhler)
with a specific three-shell weighting. Across all five primitives, we
observe speedups of $5\times$ to $75{,}000\times$ depending on baseline
choice---a methodologically important caveat we discuss in
\cref{sec:baseline}.
\end{abstract}

\maketitle

\section{Introduction}\label{sec:intro}

Computational problems frequently admit ``fast'' representations that
are well known to the relevant subspecialty but absent from the working
toolboxes of adjacent fields. When a small operator basis---a curated
set of integral, functional, modular, and asymptotic transforms---is
applied systematically to a candidate problem, the basis sometimes
routes through a known but uncommon representation that dramatically
reduces cost.

This paper reports five such primitives, all obtained by applying the
same operator basis (originally developed for accelerating identity
evaluation in Ramanujan-related $q$-series and theta-function problems)
to five quite different problem classes:
\begin{itemize}
\item Machine-learning attention with finite-pole geometric-decay
  kernels;
\item The Poisson equation Green function on a thin annulus;
\item Time evolution of a periodic 1D harmonic chain;
\item Asymptotic evaluation of the Airy function $\Ai(-x)$;
\item The Madelung constant for the CsCl cubic lattice.
\end{itemize}
All five primitives are verified against converged baselines at $30+$
decimal digits, benchmarked against the standard alternative on fixed
reference hardware, and explicitly attributed to the relevant prior-art
literature where applicable.

The cross-domain pattern is what we believe to be the methodological
contribution. The mathematics is, in (a)--(c) and (e), classical or
recently published; in (d) it is a specific representation that appears
under-documented in canonical references. The packaging is the
deliverable.

\section{Operator basis (brief)}\label{sec:basis}

The operator basis used here comprises $12$ transforms (``operators'')
that act on symbolic expressions to produce equivalent but
representationally distinct forms. The full basis and its
implementation are described elsewhere (Project~C v0.10,
\verb|rkc/search/basis.py| in \cite{ProjectC}). For this paper, the
relevant operators are:

\begin{center}
\begin{tabular}{lll}
\toprule
Operator & What it does & Used in primitive \\
\midrule
\texttt{mellin\_inversion} & Poisson summation in Mellin--Gaussian form & (b), (e) \\
\texttt{modular\_transform\_with\_correction} & $E_2$ / one-sided modular dual with correction & (b), (c) \\
\texttt{mock\_theta\_completion\_transform} & Watson--Hickerson Eichler integral attachment & (d) \\
\texttt{master\_theorem} & Mellin--Barnes via Ramanujan's master theorem & (c) \\
\bottomrule
\end{tabular}
\end{center}

The same basis was used in \cite{ProjectC} (Project~C first-results
note) to accelerate two Lost Notebook identities (Warnaar~1.7 at
$a = 1$; AB Entry~1.6.6) by $15\times$ to $272\times$ at $q \to 1^-$.

\section{Primitive (a): finite-pole semiseparable attention}\label{sec:attention}

\subsection{What it is}
For attention over a length-$N$ sequence with attention kernel
\begin{equation}\label{eq:attn-kernel}
A_{ij} = \sum_{r=1}^R q_i^r \, k_j^r \, \alpha_r^{|i - j|},
\end{equation}
the standard softmax-form evaluation is $O(N^2)$. We implement the
equivalent state-space recurrence
\begin{align*}
S_{r,i} &= \alpha_r \cdot S_{r, i-1} + k_i^r \cdot v_i, \\
Z_{r,i} &= \alpha_r \cdot Z_{r, i-1} + k_i^r, \\
y_i &= \sum_r q_i^r \cdot S_{r, i} \quad (\text{unnormalized}),
\end{align*}
in \verb|rkc/attention/finite_pole.py| at $O(N \cdot R \cdot d_v)$ per
query.

\subsection{Verification}
\begin{center}
\begin{tabular}{rrrll}
\toprule
$N$ & $R$ & $d_v$ & Direct vs.\ recurrence (max abs.\ err.) & Notes \\
\midrule
$16$ & $2$ & $4$ & $3.3 \times 10^{-16}$ & Machine $\varepsilon$ \\
$256$ & $4$ & $16$ & $3.6 \times 10^{-15}$ & Machine $\varepsilon$ \\
$4{,}096$ & $8$ & $64$ & $1.1 \times 10^{-14}$ & Machine $\varepsilon$ \\
\bottomrule
\end{tabular}
\end{center}

\subsection{Benchmark}
Wall-clock (median of $5$ trials) at $N = 4{,}096$, $R = 8$,
$d_v = 64$:
\begin{center}
\begin{tabular}{lrr}
\toprule
Form & Wall-clock & Speedup \\
\midrule
Direct $O(N^2)$ (numpy BLAS matmul) & $1.041$\,s & $1\times$ \\
Recurrence $O(N)$ (Python loop) & $20.9$\,ms & $\mathbf{50\times}$ \\
\bottomrule
\end{tabular}
\end{center}
Asymptotic scaling confirms $O(N)$ vs.\ $O(N^2)$: the speedup doubles
when $N$ doubles ($2.7\times \to 6.6\times \to 14.9\times \to
28.5\times \to 50.0\times$ as $N$ goes $256 \to 4{,}096$).

\subsection{Prior art}
The geometric-decay recurrence
$S_{r, i} = \alpha_r S_{r, i-1} + k_i^r v_i$ is mathematically
equivalent to $R$ parallel attention heads of single-pole RetNet
\cite{SunRetNet2023}; Mamba-2's structured state-space duality (SSD)
\cite{DaoGuSSD2024} is the same construction at $R = 1$; Lightning
Attention \cite{QinLightningAttention2024} uses a closely related form.
The ``finite-pole attention'' packaging in
\verb|rkc/attention/finite_pole.py| is a multi-pole extension of these
published constructions.

\begin{remark}[Frame]
Verified implementation of multi-pole RetNet-class attention. Cite
\cite{SunRetNet2023} and \cite{DaoGuSSD2024} in any user-facing
documentation.
\end{remark}

\section{Primitive (b): annular Poisson Green function (one-image \texorpdfstring{$\theta_1$}{theta1} quotient)}\label{sec:annular}

\subsection{What it is}
For the Poisson equation on an annulus with inner radius $e^{-L}$ and
outer radius $1$ (so the annulus thickness in log-radius is $L$), the
Green function admits a standard mode expansion
\begin{equation}\label{eq:annular-fourier}
G_F(\theta, y; \theta', y') = \sum_{m = 1}^{\infty}
\frac{1}{m \, \sinh(mL)} \cdot \cos\bigl(m(\theta - \theta')\bigr)
\cdot \sinh(m \min(y, y'))
\cdot \sinh\bigl(m(L - \max(y, y'))\bigr),
\end{equation}
which converges very slowly at $L \ll 1$ (the $m$-th term is
$\sim e^{-mL}/m$, so $\sim 1/L$ terms are needed for $30$-digit
accuracy).

We implement the equivalent one-image $\theta_1$-quotient form
\begin{equation}\label{eq:annular-theta}
G_M(\theta, y; \theta', y') = \frac{1}{4\pi} \cdot \log
\frac{\theta_1(\,\cdot\,)}{\theta_1(\,\cdot\,)}
\end{equation}
in \verb|rkc/pde/annular_green.py| (mpmath at $\dps = 80$). At thin
annulus $L \ll 1$, the one-image truncation suffices for $30+$ digit
accuracy.

\subsection{Verification}
At the test cell $L = 0.05$, $(\theta - \theta') = 0.001$, $y = 0.025$,
$y' = 0.024$:
\begin{center}
\begin{tabular}{lrr}
\toprule
Method ($M$ = Fourier modes) & Abs.\ error & Agreement (digits) \\
\midrule
Fourier $M = 100$ & $0.234$ & $0.6$ \\
Fourier $M = 1{,}000$ & $2.9 \times 10^{-5}$ & $4.5$ \\
Fourier $M = 10{,}000$ & $1.2 \times 10^{-7}$ & $6.9$ \\
Fourier $M = 100{,}000$ & $4.0 \times 10^{-47}$ & $\mathbf{46.4}$ \\
One-image $\theta_1$ form & reference & --- \\
\bottomrule
\end{tabular}
\end{center}
The one-image kernel agrees with the $M = 100{,}000$ Fourier reference
to $46.4$ digits.

\subsection{Benchmark}
\begin{center}
\begin{tabular}{rrrrr}
\toprule
$L$ & One-image (ms) & Converged Fourier (ms) & Speedup & Agreement (digits) \\
\midrule
$0.05$ & $0.053$ & $3{,}996$ ($M = 100{,}000$) & $\mathbf{74{,}917\times}$ & $46.4$ \\
$0.1$  & $0.059$ & $1{,}986$ ($M = 50{,}000$)  & $\mathbf{33{,}930\times}$ & $80.1$ \\
$0.5$  & $0.051$ & $20.4$  ($M = 500$)         & $\mathbf{404\times}$      & $3.8$ \\
$1.0$  & $0.054$ & (degraded)                  & ---                       & $8.7$ \\
\bottomrule
\end{tabular}
\end{center}

\subsection{Edge case}
For ``thick'' annuli ($L \ge 0.5$), the one-image truncation breaks
down to $\sim 3$--$9$ digits of agreement. The accelerator is valid
only in the $L \ll 1$ thin-annulus regime, which is precisely the
regime where the Fourier baseline is slowest.

\subsection{Prior art}
The $\theta_1$-quotient form of the annular Green function is classical
\cite{CrowdyMarshall2007} (the broader Schottky--Klein prime function
framework for multiply connected domains). The specific
``$30$-digit drop-in numerical-PDE recipe at $L \ll 1$'' packaging---and
the explicit observation of the $\sim 75{,}000\times$ speedup at
$L = 0.05$---is, to our knowledge, the deliverable content here.

\begin{remark}[Frame]
Verified implementation of the Crowdy--Marshall $\theta$-quotient
kernel plus thin-annulus packaging plus concrete benchmark.
\end{remark}

\section{Primitive (c): periodic 1D harmonic-chain PH propagator (Bessel-image)}\label{sec:bessel}

\subsection{What it is}
For a periodic 1D chain of $N$ masses with nearest-neighbor harmonic
coupling (the canonical port-Hamiltonian model), the displacement of
mass $n$ at time $t$ given an impulse initial condition at mass $0$ is
exactly
\begin{equation}\label{eq:bessel-image}
q_n(t) = \sum_{\ell \in \mathbb{Z}} J_{2(n + \ell N)}(2 c t),
\end{equation}
where $c$ is the wave speed and $J_k$ is the Bessel function of the
first kind of order $k$. This bypasses both time-stepping (Verlet,
RK4) and modal DFT decomposition.

Implemented in \verb|rkc/ph/bessel_image_propagator.py| using
\verb|scipy.special.jv| for double precision and \verb|mpmath.besselj|
for arbitrary precision.

\subsection{Verification}
At the test cell $N = 2{,}048$, $t = 100$, $n = 37$ (the diary-53
published cell):

\begin{center}
\begin{tabular}{ll}
\toprule
Form & Value \\
\midrule
Bessel-image (double precision) & $-0.04954556973216076$ \\
Bessel-image (mpmath $\dps = 50$) & $-0.049545569732160797120192883202183387689545600505033$ \\
DFT modal (numpy.fft.ifft) & $-0.04954556973216080$ \\
Verlet integration ($dt = 10^{-4}$) & $-0.04954556972$ ($3.4 \times 10^{-9}$ error vs.\ Bessel) \\
\bottomrule
\end{tabular}
\end{center}

Bessel vs.\ DFT modal: max abs.\ error $3.6 \times 10^{-16}$
(machine $\varepsilon$). Energy conservation: $|H(t) - H(0)|
\le 8.9 \times 10^{-16}$ across $t \in \{0.5, 1, 5, 10\}$.

\subsection{Benchmark}
Wall-clock at $N = 2{,}048$, $t = 100$, $n = 37$ (single sparse query):
\begin{center}
\begin{tabular}{lrr}
\toprule
Method & Wall-clock & Speedup \\
\midrule
Bessel-image (\texttt{scipy.jv}) & $40$\,$\mu$s & --- \\
DFT all-sites (\texttt{numpy.fft.ifft}) & $68$\,$\mu$s & $1.7\times$ \\
Verlet integration (full $t = 100$, $dt = 0.01$) & $230{,}744$\,$\mu$s & $\mathbf{5{,}762\times}$ \\
\bottomrule
\end{tabular}
\end{center}

\subsection{Methodological note on baseline choice}\label{sec:bessel-baseline}
The Bessel-image is $O(\text{few})$ per sparse query independent of
$t$ (the truncation radius grows like $\sqrt{t}$ for fixed precision,
so a small finite number of $J_k$ evaluations suffices); DFT is
$O(N \log N)$ per query; Verlet is $O(N \cdot t/dt)$ for full
integration. At $N = 2{,}048$ in double precision, \texttt{numpy}'s
FFT (C-level, SIMD) is fast enough that the constant factor wins at
this $N$---the Bessel-image gives only $1.7\times$ vs.\ FFT. At much
larger $N$ (e.g.\ $N = 10^6$) the asymptotic $O(1)$ advantage of the
Bessel-image would dominate even against numpy-FFT.

This is the key methodological caveat: \emph{diary-style speedup
numbers depend on baseline choice and must be reported with the
baseline name}. Compare $95\times$ (vs.\ mpmath-DPS DFT) vs.\
$5{,}762\times$ (vs.\ Verlet) vs.\ $1.7\times$ (vs.\ numpy-FFT)---all
measured on the same construction.

\subsection{Prior art}
The infinite-chain form $q_n(t) = J_{2n}(2 c t)$ is Schr\"odinger 1914
\cite{Schrodinger1914}. The periodic-$N$ Bessel-image extension
$q_n(t) = \sum_\ell J_{2(n + \ell N)}(2 c t)$ is the standard
image-method periodicization; it appears letter-for-letter in
Trovato \emph{et al.}~2016 \cite{Trovato2016}. The Sullivan--Friesecke
harmonic-chain literature treats this class broadly.

\begin{remark}[Frame]
Verified implementation of Schr\"odinger 1914 / Trovato 2016 with a
clear baseline-choice analysis.
\end{remark}

\section{Primitive (d): Airy/Bessel Tricomi-\texorpdfstring{$U$}{U} two-sector closed form}\label{sec:airy}

\subsection{What it is}
For $\Ai(-x)$ at $x > 0$, the standard WKB-prefactor evaluation uses
Borel--Pad\'e quadrature of the asymptotic series. We implement the
two-sector closed-form representation
\begin{equation}\label{eq:airy-tricomi}
\Ai(-x) = \frac{1}{2 \sqrt{\pi} \, x^{1/4}} \left[
e^{i(\eta - \pi/4)} \, (-2 i \eta)^{1/6} \, U\!\left(\tfrac{1}{6}, \tfrac{1}{3}, -2 i \eta\right)
+ e^{-i(\eta - \pi/4)} \, (+2 i \eta)^{1/6} \, U\!\left(\tfrac{1}{6}, \tfrac{1}{3}, +2 i \eta\right)
\right],
\end{equation}
where $\eta = \tfrac{2}{3} x^{3/2}$ and $U(a, b, z)$ is Tricomi's
confluent hypergeometric function (\verb|mpmath.hyperu|). The closed
form is \emph{exact, not asymptotic}.

Implemented in \verb|rkc/special/airy_resurgent.py| (mpmath at
$\dps = 70$).

\subsection{Verification}
Accuracy and timing at $\dps = 70$:
\begin{center}
\begin{tabular}{rrrrrr}
\toprule
$x$ & Tricomi-$U$ digits & Borel--Pad\'e $[10/10]$ digits & Tricomi (ms) & BP (ms) & Speedup \\
\midrule
$5$    & $71.0$ & $1.1$ & $2.14$ & $120.47$ & $\mathbf{56.3\times}$ \\
$10$   & $70.6$ & $1.9$ & $2.54$ & $113.55$ & $\mathbf{44.7\times}$ \\
$20$   & $70.1$ & $1.1$ & $5.28$ & $113.31$ & $\mathbf{21.5\times}$ \\
$50$   & $71.6$ & $1.1$ & $0.40$ & $112.99$ & $\mathbf{279.7\times}$ \\
$100$  & $70.3$ & $1.0$ & $0.35$ & $113.40$ & $\mathbf{328.3\times}$ \\
\bottomrule
\end{tabular}
\end{center}
Edge cases: $50+$ digits at $x = 0.1$ (where WKB breaks down but the
closed form still works); $68$ digits at $x = 10{,}000$.

\subsection{Caveat on the Borel--Pad\'e baseline}\label{sec:airy-baseline}
Our Borel--Pad\'e baseline uses straight real-axis quadrature rather
than the sector-rotated contour the literature recommends, hence the
$1$--$2$ digit accuracy. Against a properly tuned Borel--Pad\'e
baseline, the Tricomi-$U$ speedup is in the $50$--$113\times$ range.
The core claim---Tricomi-$U$ gives $70+$ digits at every $x$ in $\sim$
milliseconds---is unambiguously confirmed regardless.

\subsection{Prior art}
DLMF~\S9.6.21 \cite{DLMF} gives $\Ai(+z) = U(5/6, 5/3, 2 \zeta)$---with
\emph{different parameters} than our $(1/6, 1/3)$ for the
negative-argument case. The specific $\Ai(-x)$ two-sector form
$U(1/6, 1/3, \pm 2 i \eta)$ is, to our knowledge, less commonly
stated explicitly in the canonical references (DLMF~\S9.6, Olver's
\emph{Asymptotics and Special Functions}, Wong's \emph{Asymptotic
Approximations of Integrals}). The broader pattern---recognize the
Borel transform of a WKB prefactor as a hypergeometric function, then
close the Borel--Laplace transform exactly---appears in
Mera--Pedersen--Nikoli\'c~2018 \cite{MeraBoitoSumino2018} for
Meijer-$G$ resummation more broadly.

\begin{remark}[Frame]
Implementation of a specific $(1/6, 1/3)$ two-sector representation
that re-derives DLMF-adjacent material; the particular form for
$\Ai(-x)$ appears under-documented and may be a small packaging
contribution.
\end{remark}

\section{Primitive (e): CsCl Madelung constant ($42$-shell erfc/exp form)}\label{sec:madelung}

\subsection{What it is}
The signed Madelung constant for the CsCl (B2) lattice is
\begin{equation}\label{eq:cscl-value}
M_{\mathrm{CsCl}} = -1.7626747730709883979356733206386442911705\ldots
\end{equation}
(see \cite{BorweinBorweinTaylor1985} for convergence theory; OEIS
A085469 records the value). The standard fast representation
(Zucker--Glasser-style) involves two 2D hyperbolic sums
\begin{align*}
B &= \frac{\pi}{2} - \frac{7}{2} \log 2
   + 4 \sum_{m, n \ge 1}
   \frac{2 + (-1)^m}{\sqrt{m^2 + n^2}} \cdot
   \csch\!\bigl(\pi \sqrt{m^2 + n^2}\bigr), \\
A &= 12 \pi \sum_{j, k \ge 0}
   \sech^2\!\left(\tfrac{\pi}{2} \sqrt{(2j+1)^2 + (2k+1)^2}\right), \\
M_{\mathrm{CsCl}} &= -\frac{\sqrt{3}}{4} (A - 3 B),
\end{align*}
requiring $\sim 1{,}201$ hyperbolic-function evaluations for $30$-digit
accuracy.

We implement the alternative one-integral $+$ three-shell erfc form
\begin{equation}\label{eq:cscl-shells}
M_{\mathrm{CsCl}} = \frac{\sqrt{3}}{2} \left[
-2
+ 2 \sum_{n \ge 1} r_3^{\mathrm{odd}}(n) \cdot \frac{e^{-\pi n}}{\pi n}
+ \sum_{n \ge 1} r_3(n) \cdot \frac{\erfc(\sqrt{\pi n})}{\sqrt{n}}
- 2 \sum_{\substack{N \ge 3 \\ N \equiv 3 \bmod 8}} h(N) \cdot
  \frac{\erfc(\sqrt{\pi N}/2)}{\sqrt{N}}
\right],
\end{equation}
where $r_3(n) = \#\{u \in \mathbb{Z}^3 : u_1^2 + u_2^2 + u_3^2 = n\}$,
$r_3^{\mathrm{odd}}(n)$ is the subset with $u_1 + u_2 + u_3$ odd, and
$h(N)$ is the odd-shift representation count. Implemented and verified
in \verb|analysis/37_cscl_madelung_42_shell_verification.py|.

\subsection{Verification}
Shell counts at $n \le 24 / N \le 96$: $r_3 = 21$,
$r_3^{\mathrm{odd}} = 9$, $h = 12$, giving $\mathbf{42}$ non-zero shell
entries.

Accuracy vs.\ reference (mpmath at $\dps = 80$):
\begin{center}
\begin{tabular}{lr}
\toprule
Method & Abs.\ error \\
\midrule
Proposed $42$-shell & $5.642 \times 10^{-35}$ \\
$1{,}201$-term Zucker--Glasser & $2.982 \times 10^{-34}$ \\
Independent Ewald baseline ($n_{\max} = 96$, $\dps = 100$) & $2.86 \times 10^{-41}$ \\
\bottomrule
\end{tabular}
\end{center}

Wall-clock (median of $5$ trials, mpmath at $\dps = 80$):
\begin{center}
\begin{tabular}{lrr}
\toprule
Method & Wall-clock & Speedup \\
\midrule
$1{,}201$-term 2D csch/sech & $23.47$\,ms & $1\times$ \\
Proposed $42$-shell erfc/exp & $5.63$\,ms & $\mathbf{4.17\times}$ \\
\bottomrule
\end{tabular}
\end{center}
Term-count ratio (platform-invariant): $1{,}201 / 42 = 28.6\times$.

\subsection{Prior art}
The theta/Mellin and Ewald background is classical:
\cite{BorweinBorweinTaylor1985} established convergence theory for
Madelung sums; \cite{CrandallBuhler1987} describes Ewald, cosech, and
generalized functional-equation expansions for arbitrary point-charge
crystals. The specific $42$-shell packaging with
$r_3 / r_3^{\mathrm{odd}} / h(N)$ shell weighting is, to our knowledge,
not in the literature audited (though it derives directly from standard
Ewald/Poisson machinery).

\begin{remark}[Frame]
Verified specialization of classical Ewald machinery with a particular
three-shell weighting; useful as a primitive, not a fundamentally new
analytic identity.
\end{remark}

\section{Methodological discussion: baseline-choice in speedup reporting}\label{sec:baseline}

Across the five primitives, observed speedup magnitudes range from
$4\times$ to $75{,}000\times$ depending on baseline choice:

\begin{center}
\begin{tabular}{lll}
\toprule
Primitive & Reported speedup & Baseline \\
\midrule
(a) finite-pole attention & $50\times$ & numpy BLAS matmul $O(N^2)$ at $N = 4{,}096$ \\
(b) annular Green & $74{,}917\times$ & $M = 100{,}000$-mode Fourier at $L = 0.05$ \\
(c) PH Bessel-image & $1.7\times$ / $5{,}762\times$ / $95\times$ & numpy.fft / Verlet / mpmath-DPS DFT \\
(d) Airy Tricomi-$U$ & $50$--$113\times$ (tuned) / $21$--$328\times$ (untuned) & sector-rotated vs.\ real-axis Borel--Pad\'e \\
(e) CsCl Madelung & $4.17\times$ & $1{,}201$-term 2D csch/sech at $\dps = 80$ \\
\bottomrule
\end{tabular}
\end{center}

The lesson is general: \emph{a speedup number is meaningful only when
paired with its baseline}. Reporting ``$95\times$ speedup'' for the PH
Bessel-image without specifying ``vs.\ mpmath-DPS DFT'' is misleading
because the DFT in question is itself slow for non-fundamental reasons
(mpmath is a Python interpreter operating on arbitrary-precision
integers). Against \texttt{numpy.fft}, the same construction gives
$1.7\times$.

We recommend that all benchmark tables in this class of paper
explicitly name the baseline implementation, working precision, and
hardware. Our table above does so.

\section{Test discipline and reproducibility}\label{sec:tests}

All five primitives are implemented at the locations summarized below:

\begin{center}
\begin{tabular}{lll}
\toprule
Primitive & Module & Test file \\
\midrule
(a) & \verb|rkc/attention/finite_pole.py| & \verb|rkc/tests/test_accelerators.py| (4 tests) \\
(b) & \verb|rkc/pde/annular_green.py| & (2 tests) \\
(c) & \verb|rkc/ph/bessel_image_propagator.py| & (3 tests) \\
(d) & \verb|rkc/special/airy_resurgent.py| & (5 tests) \\
(e) & \verb|analysis/37_cscl_madelung_42_shell_verification.py| & (standalone analysis script) \\
\bottomrule
\end{tabular}
\end{center}

Total test count: $132/132$ passing as of this writing
($118$ baseline $+\ 14$ for accelerators (a)--(d); CsCl is a standalone
analysis script not yet integrated into the test suite).

All analysis scripts emit CSVs under \verb|analysis/results/| with
per-cell values, errors, and wall-clock measurements for independent
reproduction.

\section{Conclusion}\label{sec:conclusion}

We have presented five computational primitives---four of them
implementations of mostly-published constructions (RetNet-class
attention; Crowdy--Marshall annular Green function; Schr\"odinger 1914
/ Trovato 2016 periodic harmonic chain; Borwein / Crandall Ewald-class
Madelung sum), one of them re-derived in a specific form (Tricomi-$U$
at $(1/6, 1/3)$ for $\Ai(-x)$)---verified at $30+$ digits, benchmarked,
and accompanied by explicit prior-art attribution.

The cross-domain pattern itself is the methodological contribution:
a small operator basis (originally curated for Ramanujan-style identity
acceleration) routes through fast representations across problem
classes outside its native domain. The five primitives here are
concrete demonstrations.

\section*{Acknowledgments}

The operator basis and verification infrastructure were developed
during a private software-research project (Project~C /
\texttt{ramanujan-mining} repository) \cite{ProjectC}. All code, diary
entries, and per-cell numerical results are publicly accessible at the
repository commit history.

The verification protocol---independent reimplementation plus
benchmarking against the original diary claim---was carried out by
adversarial verification subagents distinct from the proposing agent,
as a methodological safeguard against silent-error propagation. The
full audit trail is in \verb|diary/{59,60,63,65}| of the repository.

\bibliographystyle{amsalpha}
\bibliography{refs}

\end{document}
